\section{الفصل~\ref{sec:dofaalt}. نظريات أخرى في \nt{DOPA}}

كل واحدة من الصور الموسَّعة تستند إلى قياساتها المباشرة للدوبامين (تشتت نقاط الانقلاب، والاستجابات للمزعج والجديد، والارتفاع البطيء، والاستجابات لتغيّر الطعم)، لكن الصيغ التي تفسرها نظريات، وبين اثنتين منها تتناقض التجارب حتى الآن.

{\footnotesize
\setlength{\tabcolsep}{3pt}\renewcommand{\arraystretch}{1.15}
\begin{longtable}{>{\raggedright}p{0.5cm}>{\raggedright}p{4.0cm}>{\raggedright}p{5.6cm}>{\raggedright}p{2.3cm}>{\raggedright\arraybackslash}p{1.7cm}}
\toprule
م & القضية & التجربة وما قيس & المصدر & الحالة \\
\midrule\endfirsthead
\toprule
م & القضية & التجربة وما قيس & المصدر & الحالة \\
\midrule\endhead
1 & القاعدة غير المتناظرة~\eqref{eq:asymmetric} تتقارب إلى النقطة~\eqref{eq:expectile}؛ عند $\kappa=1/2$ هي المتوسط، ولقطرة باحتمال $1/2$ تكون $V=\kappa$ (الشكل~\ref{fig:expectiles}) & النقطة~\eqref{eq:expectile} هي ”التوقّعي“ (\foreignlanguage{english}{expectile})، أي حل مسألة المربعات الصغرى بأوزان مختلفة للانحرافات الموجبة والسالبة؛ ولمثال الفصل الاشتقاق في الفصل نفسه. & \cite{NeweyPowell1987}؛ اشتقاق في الفصل & نظرية \\
2 & لعصبونات \nt{DOPA} المختلفة نقاط انقلاب مختلفة، مرتبة بحسب اللاتناظر (القضية~\ref{prop:expectile}) & فئران، عصبونات \nt{DOPA} في السقيفة البطنية موسومة بصريًا جينيًا، ومكافآت مختلفة الحجم: النقطة التي تتحول فيها الذروة إلى هبوط تختلف من عصبون إلى آخر؛ وهي أعلى عند العصبونات ذات اللاتناظر الأكبر في الاستجابات؛ ومن استجابات المجموعة يُستعاد شكل توزيع المكافآت. أما أن هذا هو الوظيفة الرئيسية للتشتت ففرضية. & \cite{Dabney2020} & مُقاس \\
3 & جزء من عصبونات \nt{DOPA} يستجيب للمزعج بذروة & قرود: بعض عصبونات \nt{DOPA} يُستثار بالإشارة المنبئة بالمكافأة ويُثبَّط بالإشارة المنبئة بنفخة هواء في العين، وبعضها يُستثار بكلتيهما؛ والمجموعتان في أجزاء مختلفة من الدماغ المتوسط. & \cite{Matsumoto2009, BrombergMartin2010} & مُقاس \\
4 & القيمة المطلقة للخطأ أو الجِدّة تحدد معدل التعلّم & لا توجد في الأعمال المستشهد بها تجربة مباشرة تغيّر فيها إشارة الأهمية عند عصبونات \nt{DOPA} معدل التعلّم؛ والمراجعة تقترح دور إشارة ”انتبه، هذا مهم“. & \cite{BrombergMartin2010} & فرضية \\
5 & سلك مستقل لمحور الخطر & فئران: عصبونات \nt{DOPA} الذاهبة إلى الطرف الخلفي من المخطط تستجيب للمنبّهات الجديدة والقوية لا للمكافأة؛ وإتلافها يضعف تجنّب الشيء الجديد. & \cite{Menegas2018} & مُقاس \\
6 & زمن الفعل الأمثل $\tau_a^*=\sqrt{C/\bar R}$~\eqref{eq:vigor} & أصغر قيمة للمجموع $C/\tau_a+\bar R\tau_a$؛ والاشتقاق في الفصل. في النموذج الأصلي يحدد سرعةَ الأفعال ثمنُ الوقت، وهو يساوي متوسط المكافأة. & \cite{Niv2007}؛ اشتقاق في الفصل & نظرية \\
7 & المستوى البطيء للدوبامين يحمل $\bar R$ ويحدد سرعة الأفعال (القضية~\ref{prop:multiplex}) & جرذان، غسيل مجهري وقياس فولتية في النواة المتكئة: مستوى الدوبامين يتبع من دقيقة إلى دقيقة تواترَ المكافآت وسرعةَ الأفعال. وعند البشر يقوّي \foreignlanguage{english}{L-DOPA} اعتماد زمن الاستجابة على متوسط تواتر المكافآت. أما أن المستوى البطيء يساوي $\bar R$ تحديدًا فغير مبرهن. & \cite{Hamid2016, Beierholm2013vigor} & فرضية \\
8 & المستوى البطيء يتألف من مصدرين: التحرير يتغير عند المخارج دون تغيّر تردد النبضات، لكن الارتفاع البطيء موجود في التردد نفسه أيضًا & جرذان، مهمة واحدة: التحرير في النواة المتكئة يتبع التوقع المتغير للمكافأة، أما تردد نبضات عصبونات \nt{DOPA} المعرَّفة في السقيفة البطنية فلا. فئران في ممر افتراضي (السطر~10): الارتفاع التدريجي عند الاقتراب من المكافأة موجود في نبضات عصبونات \nt{DOPA} المعرَّفة أيضًا. & \cite{Mohebi2019, Kim2020} & مُقاس \\
9 & يرتفع الدوبامين تدريجيًا ما دام الحيوان يقترب من مكافأة بعيدة & جرذان في متاهة، قياس الفولتية في المخطط: يتزايد تركيز الدوبامين خلال ثوانٍ مع الاقتراب من الهدف، ويعتمد الانحدار على المسافة وعلى حجم المكافأة. & \cite{Howe2013ramp, Hamid2016} & مُقاس \\
10 & الارتفاع هو الخطأ $\delta$ مع تدقيق التقدير~\eqref{eq:ramp}؛ والقفزة في الممر تعطي ذروة (الشكل~\ref{fig:ramp}) & الصيغة والعدد $\delta=0.75\,V$ في الثانية اشتقاق في الفصل. التجربة: فئران في ممر افتراضي، ونقل لحظي إلى نقطة أقرب من الهدف؛ نبضات أجسام الخلايا، والكالسيوم في الأجسام والمخارج، والتركيز في المخطط تستجيب كخطأ لا كتوقع $V$. أي التفسيرين صحيح على كل المخارج، مسألة قيد النقاش. & \cite{Kim2020} & مُقاس \\
11 & النظر إلى الوراء: يخبر الدوبامين بمقياس الصلة السببية~\eqref{eq:retro} & فئران، تحرير الدوبامين في النواة المتكئة في تجارب تتنبأ فيها هذه النظرية والخطأ~\eqref{eq:ctd} بنتائج مختلفة: في عدد من التجارب تبع التحرير الأولى. & \cite{Jeong2022} & فرضية \\
12 & تنزاح استجابة الدوبامين أثناء التعلّم تدريجيًا من المكافأة إلى الإشارة المنبئة & فئران، إشارة مخارج عصبونات \nt{DOPA} في المخطط البطني على مدى التعلّم: تنتقل الذروة تدريجيًا من لحظة المكافأة إلى لحظة الإشارة المنبئة، كما يقتضي التعلّم وفق~\eqref{eq:ctd}. & \cite{Amo2022} & مُقاس \\
13 & عصبونات \nt{DOPA} تستجيب لتغيّر طعم المكافأة مع ثبات القيمة & جرذان، تسجيل عصبونات السقيفة البطنية: استبدال عصير بآخر مساوٍ له في القيمة يستثير استجابة، مع أن الخطأ~\eqref{eq:ctd} يساوي صفرًا. & \cite{Takahashi2017} & مُقاس \\
14 & الذرى الاصطناعية تعلّم الصلة بين منبّهين محايدين & جرذان، تجربة إقران مسبق لمنبّهين دون مكافأة: التشغيل البصري الجيني لعصبونات \nt{DOPA} عند الانتقال من الأول إلى الثاني ينشئ الصلة، وإيقافها يعيق تعلّمها. & \cite{Sharpe2017} & مُقاس \\
15 & متجه الأخطاء بحسب السمات~\eqref{eq:vectortd} & نظرية خطأ التنبؤ بالسمات؛ ولا يوجد اختبار مباشر للصيغة المتجهية. & \cite{Gardner2018} & فرضية \\
16 & عصبونات \nt{DOPA} المختلفة في المهمة الواحدة تحمل مقادير مختلفة & فئران في ممر افتراضي، تصوير ثنائي الفوتون لعصبونات \nt{DOPA}: بعضها مرتبط بالمكافأة، وبعضها بالسرعة والانعطاف، وبعضها بالإشارات المنبئة ودقة الاختيار. & \cite{Engelhard2019} & مُقاس \\
17 & حزمة أسلاك بدل سلك (القضية~\ref{prop:bundle}) & خلاصة السطور~\babelsublr{2--16}: كل تفكيك مقيس على حدة؛ ولا توجد تجربة تبيّن صورة موحدة تكون فيها كلها أجزاءً من إشارة واحدة. & --- & فرضية \\
\bottomrule
\end{longtable}}
